\section{The global harmonic filtration}\label{sec:global}

We now filter the entire spherical-coset poset.  Throughout the construction
assume that $S\ne\varnothing$; the empty system will be included in the final
theorem.  Recall that $\mathcal D$ consists of the right cosets $A_Tx$, with
$T$ spherical, ordered by inclusion with nested types.  All realizations of
posets have the CW topology.

\subsection{Harmonic vectors and a well-ordered height}

For a spherical subset $T\subseteq S$ and a vector
$v\in\R^{\widehat S}$, let $\mathsf H_T(v)$ be the vector that agrees with
$v$ outside $T$ and is harmonic on $T$.  Thus
\begin{equation}\label{eq:global-harmonization}
 \mathsf H_T(v)_{\widehat S\setminus T}=v_{\widehat S\setminus T},
 \qquad
 \mathsf H_T(v)_T
   =C_T^{-1}\left(\sum_{j\notin T}b_{ij}v_j\right)_{i\in T}.
\end{equation}
For $T=\varnothing$, this is the identity.  Nonnegativity of $C_T^{-1}$
shows that $\mathsf H_T$ preserves nonnegative vectors.

For $P=A_Tx\in\mathcal D$ and every integer $k$, put
\begin{equation}\label{eq:global-coset-vector}
 p^k(P)=\mathsf H_T\bigl(m^k(x)\bigr).
\end{equation}
Operations in a color of $T$ change the layer multiplicities only in that
color.  Therefore all $m_i^k(x)$ with $i\notin T$ are constant on $A_Tx$,
and \eqref{eq:global-coset-vector} is independent of its representative.
The vectors $p^k(P)$ are nonnegative and have finite support in $k$.
The frame coordinate satisfies
\begin{equation}\label{eq:global-frame}
 p_o^k(P)=m_o^k(x)=
 \begin{cases}1,&k=0,\\0,&k\ne0.\end{cases}
\end{equation}

For example, the identity representative has $m^0(1)=\varepsilon_o$
and all other layers zero, because its family consists of the frame
projectives themselves.  On the coset $A_T$ this gives
$p_T^0(A_T)=C_T^{-1}\mathbf1$, with frame coordinate $1$ and all
other exterior coordinates zero.  In rank one the interior coordinate
is $1/2$.  This is the simplest instance of replacing generator counts
by harmonic coset data.

We shall repeatedly use two harmonic comparison formulas.  If
$P=A_Tx\subseteq A_Ux=P'$ with $U$ spherical, then the difference
$p^k(P')-p^k(P)$ is supported on $U$, and
\begin{equation}\label{eq:global-upper-change}
 \bigl(p^k(P')-p^k(P)\bigr)_U
   =-C_U^{-1}\bigl(Cp^k(P)\bigr)_U.
\end{equation}
Indeed the two vectors agree outside $U$ and the new vector has zero
$C$-rows on $U$.

For a lower coset $Q=A_Rz\subseteq P=A_Tx$, set $V=T\setminus R$.
The two vectors agree outside $T$.  Writing
$\delta_i=m_i^k(z)-p_i^k(P)$ for $i\in V$, their difference
$\Delta=p^k(Q)-p^k(P)$ satisfies
\begin{equation}\label{eq:global-lower-change}
 \Delta_V=\delta_V,
 \qquad
 \Delta_R=C_R^{-1}
       \left(\sum_{j\in V}b_{ij}\delta_j\right)_{i\in R}.
\end{equation}
Both formulas also hold when one of the displayed types is empty.
In particular, componentwise nonnegative or nonpositive changes on the
open coordinates extend with the same sign.  A nonzero change on an open
coordinate remains nonzero in the full vector.

Fix an ordering of $\widehat S$.  For a nonnegative vector, its
positive-support components mean the connected components of the graph
on its positive coordinates, using the edges with $b_{ij}>0$.  Boundaries
of these components are taken in $\widehat S$, and may contain the frame.

The comparison formulas suggest ordering cosets by their harmonic vectors,
from higher nonnegative layers down, but equal vectors still require a choice.  Suppose
that adjoining a color with positive coordinate leaves the vector at the
comparison level unchanged.  Our ordinary preference will put the larger
type first, by assigning $-1$ to each positive coordinate belonging to the
type.  Zero coordinates contribute nothing to this tie; they may leave the
comparison to a lower level.

We reverse this preference on a positive-support component $I\subseteq S$
at a level $k>0$ when $p^k(P)$ is harmonic on $I$ and all higher vectors
vanish on $I\cup\partial I$.  The vector on such a component is a positive
null vector for its Cartan matrix.  The upper-link proof must exclude this
singular alternative when it retracts the earlier upper part to extensions
contained in one spherical type; this is why the preference cannot be uniform.  Assigning
$+1$ on this component instead puts those tied additions later.  The same
choice puts removal of an equal positive coordinate earlier, so the lower
comparison must allow equality with the harmonic bound.  This is the role
of the special closed sublevels in
Proposition~\ref{prop:residue-sublevels}, whose nonemptiness comes from the
isolated-layer obstruction.  The upper and lower verifications appear in
Lemmas~\ref{lem:global-active-spherical} and~\ref{lem:global-residue-blocks}.
The following definition encodes both preferences in the type-dependent
integer at each level.

\begin{definition}\label{def:global-sign-height}
For $P=A_Tx$, $k\ge0$, and $i\in\widehat S$, define a sign
$\eta_i^k(P)$ as follows.  It is zero if $p_i^k(P)=0$.  If $p_i^k(P)>0$,
let $I$ be its positive-support component at level $k$.  Set
$\eta_i^k(P)=+1$ precisely when
\begin{equation}\label{eq:global-null-sign}
 k>0,\qquad I\subseteq S,\qquad
 \bigl(Cp^k(P)\bigr)_I=0,\qquad
 p_j^d(P)=0
 \quad(j\in I\cup\partial I,\ d>k).
\end{equation}
In every other positive-coordinate case set $\eta_i^k(P)=-1$.
Define the level block
\begin{equation}\label{eq:global-height-block}
 \mathsf b_k(P)=
 \left(
    \sum_{i\in\widehat S}p_i^k(P),\
    (p_i^k(P))_{i\in\widehat S},\
    \sum_{i\in T}\eta_i^k(P)
 \right).
\end{equation}
Within a block use lexicographic order: first the sum, then the vector in
the fixed coordinate order, then the integer.  Compare the finite-support
sequences $(\mathsf b_k(P))_{k\ge0}$ at their highest differing level,
with the smaller block earlier.  Finally, choose a well-order of each set
of cosets with identical block sequences.
\end{definition}

Thus the sign comparison at level $k$ precedes every comparison at a
lower level.  The signs are functions only of the vectors at level $k$
and above; if those vectors coincide for two cosets, the same signs are
used at that level.

\begin{lemma}\label{lem:global-well-order}
Definition~\ref{def:global-sign-height} gives a well-order of the vertices
of $\mathcal D$.  Distinct comparable cosets have different block
sequences, so the final tie refinement does not affect either of their
earlier links.
\end{lemma}

\begin{proof}
The weighted layer counts of Section~\ref{sec:height-data} have only
finitely many possible values in a bounded interval, uniformly over the
level.  Indeed, for any fixed unfolded label the count is a nonnegative
integer combination of finitely many fixed positive weights, up to a
fixed positive proportionality factor.  A bound on the count bounds each
integer coefficient.  This proves the asserted finiteness without
requiring the real weights to lie in a lattice.

If $\sum_i p_i^k(P)\le B$, every exterior coordinate
$m_i^k(x)=p_i^k(P)$ for $i\notin T$ lies in $[0,B]$.  There are therefore
only finitely many exterior vectors.  For a fixed type $T$,
\eqref{eq:global-harmonization} determines the interior vector uniquely.
There are only finitely many spherical types.  Consequently there are
only finitely many possible primary vectors with sum at most $B$,
uniformly over $P$ and $k$.  The secondary integer in
\eqref{eq:global-height-block} lies between $-|S|$ and $|S|$.

The set of possible level blocks thus has only finitely many predecessors
below any one of its members.  Adjoin the zero block if necessary.  It is
the least block: a zero primary vector has all signs zero, whereas every
nonzero primary vector has positive sum.  Assign to each block its rank
$n(\mathsf b)\in\Zbb_{\ge0}$ in this order, with $n(0)=0$.  For a
finite-support sequence, the ordinal
\[
 \omega^N n(\mathsf b_N)+\omega^{N-1}n(\mathsf b_{N-1})
       +\cdots+\omega n(\mathsf b_1)+n(\mathsf b_0)
       <\omega^\omega
\]
has exactly the prescribed highest-level-first comparison.  The block
sequences are therefore well ordered.  A well-ordered refinement of each
fiber is again a well-order.  Such a refinement may also be obtained from
an enumeration of words in the finitely generated group and the finitely
many types.

It remains to examine comparable ties.  For every $i\in T$, the frame
and \eqref{eq:global-harmonization} give
\begin{equation}\label{eq:global-level-zero-positive}
 p_T^0(P)\ge C_T^{-1}\mathbf1>0.
\end{equation}
All signs on positive coordinates at level zero are $-1$, so the
secondary value at that level is exactly $-|T|$.  Suppose
$A_Tx\subseteq A_Ux$ and all nonnegative primary vectors agree.
If $T\ne U$, their level-zero secondary values differ.  If $T=U$,
the two nested cosets are equal.  Hence distinct comparable vertices are
never fully tied.  Negative levels may distinguish further incomparable
cosets, but their omission cannot change this conclusion or the
well-ordering argument.
\end{proof}

Write $Q\prec P$ for the resulting order of vertices.  We must show that
the earlier strict lower and strict upper parts at every $P$ are each
empty or contractible.

\subsection{The earlier upper part}

Fix $P=A_Tx$.  For each $i\in S\setminus T$, let
\begin{equation}\label{eq:global-incidence-level}
 H_i=\max\{k\ge0:
       p_i^k(P)>0\text{ or }p_j^k(P)>0
       \text{ for some }j\text{ with }b_{ij}>0\}.
\end{equation}
This maximum exists: the vectors have finite level support, and the frame
is a positive neighbor at level zero.  Call $i$ \emph{active} when
\begin{equation}\label{eq:global-active}
 (Cp^{H_i}(P))_i>0,
 \qquad\text{or}\qquad
 (Cp^{H_i}(P))_i=0\ \text{ and }\ \eta_i^{H_i}(P)=-1.
\end{equation}
Let $\operatorname{Act}(P)$ be the set of active colors.  If the row in
\eqref{eq:global-active} is zero, the coordinate $p_i^{H_i}(P)$ is
positive: otherwise the zero row would force all its neighboring
coordinates to vanish, contrary to the definition of $H_i$.

\begin{lemma}\label{lem:global-upper-comparison}
Every earlier upper extension of $P$ adds an active color.  Every
spherical upper extension adding a nonempty set of active colors and no
other colors is earlier than $P$.
\end{lemma}

\begin{proof}
Let $P'=A_Ux$ with $T\subsetneq U$ spherical, and put
$H=\max_{i\in U\setminus T}H_i$.  At levels above $H$, the coordinates
at every added color and all of its neighbors vanish.  The $T$-rows of
$Cp^k(P)$ vanish by harmonicity, and the added rows vanish by this
incidence condition.  Formula~\eqref{eq:global-upper-change} shows that
the vectors agree above $H$.  All added signs are zero there, so the
secondary values agree as well.

At level $H$, an added color with $H_i<H$ has zero coordinate, zero row,
and zero sign.  Suppose first that no added color is active.  An added
color with $H_i=H$ has row at most zero.  Thus all rows on $U$ are
nonpositive, and \eqref{eq:global-upper-change} gives a nonnegative
change of vector.  If any such row is negative, the positive diagonal of
$C_U^{-1}$ makes this change nonzero, and the sum increases strictly.
If all these rows vanish, the vectors agree at $H$.  Each added color
with $H_i=H$ then has positive coordinate and sign $+1$.  There is at
least one such color.  Since the signs for the two cosets agree whenever
their vectors at and above $H$ agree, the secondary value increases
strictly.  In either case $P'$ is later.

If all added colors are active, the same argument reverses the
inequalities.  The rows at level $H$ are nonnegative, so the vector change
is nonpositive.  A nonzero row decreases the sum strictly.  If all rows
are zero, every added color with $H_i=H$ has sign $-1$, and their
nonempty contribution decreases the secondary value strictly.  Hence
$P'\prec P$.
\end{proof}

\begin{lemma}\label{lem:global-active-spherical}
The subset $T\cup\operatorname{Act}(P)$ is spherical.
\end{lemma}

\begin{proof}
Let $G$ be a connected component of this subset, and choose the highest
level $k\ge0$ having a positive coordinate on $G$ or on one of its
neighbors.  Such a level exists by the frame.  For a vertex of $T$ the
full $C$-row is zero.  At an active vertex whose incidence maximum is
less than $k$, its coordinate and neighboring coordinates are all zero
at $k$.  At an active vertex with incidence maximum $k$, the full row is
nonnegative by \eqref{eq:global-active}.  Therefore
\begin{equation}\label{eq:global-full-row}
 (Cp^k(P))_G\ge0,
 \qquad C_Gp_G^k(P)\ge0.
\end{equation}

The vector $p_G^k(P)$ is nonzero.  Otherwise a positive incident exterior
coordinate, whose existence would be forced by the choice of $k$, would
give a negative full row at its adjacent vertex of $G$.  Moreover, a
zero coordinate of $p_G^k(P)$ cannot neighbor a positive coordinate in
$G$, since its $C_G$-row would be negative.  Connectedness implies
$p_G^k(P)>0$.

Apply Lemma~\ref{lem:cartan-positive-null}.  Either $C_G$ is positive
definite, or $C_Gp_G^k(P)=0$.  In the latter case the full rows equal
minus the exterior incidence:
\[
 (Cp^k(P))_i=-\sum_{j\notin G}b_{ij}p_j^k(P)
 \quad(i\in G).
\]
Together with \eqref{eq:global-full-row}, this forces every displayed
sum to vanish.  There is no positive exterior neighbor, and the full
rows are zero.  Since $p_G^k(P)>0$, the set $G$ is exactly one
positive-support component of $p^k(P)$.

The level $k$ cannot be zero: the frame lies outside $G\subseteq S$,
has positive level-zero coordinate, and neighbors every vertex of $G$.
By the maximality of $k$, all higher coordinates vanish on
$G\cup\partial G$.  Thus $G$ satisfies every condition in
\eqref{eq:global-null-sign}, and all its signs at level $k$ are $+1$.
Each active vertex in $G$ would now have $H_i=k$, positive coordinate,
and zero full row, contradicting its active condition.  Hence there are
no active vertices in $G$, so $G\subseteq T$.  This contradicts the
assumption that $C_G$ is not positive definite, because $T$ is
spherical.  We conclude that every component $G$ has positive-definite
Cartan matrix, proving the assertion.
\end{proof}

\begin{proposition}\label{prop:global-upper-link}
The earlier strict upper part of $P$ is empty or contractible.
\end{proposition}

\begin{proof}
If $\operatorname{Act}(P)=\varnothing$, there is no earlier upper
extension by Lemma~\ref{lem:global-upper-comparison}.  Otherwise put
$K=T\cup\operatorname{Act}(P)$, which is spherical by
Lemma~\ref{lem:global-active-spherical}.  On the earlier upper part the
map
\[
 A_Ux\longmapsto A_{U\cap K}x
\]
is an order map pointwise below the identity.  Its image is still in the
earlier strict upper part: $U\cap K$ contains an added active color and
adds only active colors.  The map is a retraction onto the upper
extensions adding only active colors, and that image has largest element
$A_Kx$.  The order homotopy observation in Section~\ref{sec:residues}
proves contractibility.
\end{proof}

\subsection{The earlier lower part}

For $i\in T$, define
\begin{equation}\label{eq:global-prescribed-top}
 q_i=\max\{k\ge0:p_i^k(P)>0\}.
\end{equation}
The maximum exists and is nonnegative by
\eqref{eq:global-level-zero-positive}.  On a connected component $B$ of
$T$, the vector $p_B^k(P)$ is either zero everywhere or positive
everywhere: apply the strict positivity of $C_B^{-1}$ to the nonnegative
exterior forcing vector.  Consequently $q_i$ is constant on $B$, and
the signs $\eta_i^{q_i}(P)$ are also constant there, because $B$ lies in
one positive-support component at that level.

\begin{definition}\label{def:global-eligible}
For a lower coset $Q=A_Rz\subsetneq P$, call an open color
$i\in T\setminus R$ \emph{eligible} if
\begin{equation}\label{eq:global-eligibility}
 \begin{gathered}
 m_i^d(z)=0\quad\text{for every }d>q_i,\\
 m_i^{q_i}(z)<p_i^{q_i}(P),
 \quad\text{or}\quad
 m_i^{q_i}(z)=p_i^{q_i}(P)
       \ \text{ and }\ \eta_i^{q_i}(P)=+1.
 \end{gathered}
\end{equation}
The first line is required in both alternatives of the second line.
Write $E(Q)$ for the eligible open colors.  This definition is
independent of $z$, since all its tested colors lie outside $R$.
\end{definition}

\begin{lemma}\label{lem:global-lower-comparison}
Every earlier lower coset has an eligible open color.  A lower coset all
of whose open colors are eligible is earlier than $P$.  The earlier
lower part retracts by an order map pointwise above the identity onto
the poset of these all-eligible lower cosets.
\end{lemma}

\begin{proof}
Let $Q=A_Rz\subsetneq P$ and $V=T\setminus R$.  Test the comparison at
the level
\[
 H=\max\left(
    \{q_i:i\in V\}\ \cup\
    \{k\ge0:m_i^k(z)>0\text{ for some }i\in V\}
   \right).
\]
This is a finite maximum over a nonempty set.  Above it all open
coordinates of both vectors vanish.  Formula~\eqref{eq:global-lower-change}
gives equality of the full vectors there, and the removed colors have
zero signs, so the secondary values also agree.

Suppose first that $E(Q)=\varnothing$.  At level $H$, a color with
$q_i<H$ has prescribed value zero and actual value nonnegative.  A color
with $q_i=H$ has no actual data above $q_i$ by the choice of $H$;
ineligibility therefore implies $m_i^H(z)\ge p_i^H(P)$.  Every open
difference at $H$ is nonnegative.  If one is positive, the full vector
has a strictly larger sum by \eqref{eq:global-lower-change}.

If all open differences vanish, no color with $q_i<H$ has actual data
at $H$.  In particular $H$ is the prescribed maximum of at least one
open color; it could not have arisen solely from actual data above a
prescribed maximum.  The open colors with $q_i<H$ contribute zero signs
at $H$.  Each with $q_i=H$ has equal positive actual and prescribed
coordinates, and ineligibility forces sign $-1$.  The full vectors now
agree at and above $H$, so the signs are common to the two cosets.  The
change in the secondary value is
\[
 \sum_{i\in R}\eta_i^H(Q)-\sum_{i\in T}\eta_i^H(P)
       =-\sum_{i\in V}\eta_i^H(P)>0.
\]
Thus $Q$ is later in this case as well.  This proves the first assertion
by contraposition, including the possibility of actual data above the
prescribed maximum on a different open component.

If every open color is eligible, the actual caps imply
$H=\max_{i\in V}q_i$.  Open colors with smaller $q_i$ have zero actual
and prescribed values at $H$; all other open differences are
nonpositive.  A strict difference decreases the sum.  If every
difference vanishes, the nonempty set of open colors with $q_i=H$ has
sign $+1$, and its removal strictly decreases the secondary value.
Hence $Q\prec P$.

For an earlier lower coset define
\begin{equation}\label{eq:global-closure}
 \rho(Q)=A_{T\setminus E(Q)}z.
\end{equation}
Since $E(Q)$ is nonempty and contained in $T\setminus R$, this is a
proper lower coset above $Q$.  Its open colors are exactly $E(Q)$ and
remain eligible, so it is earlier.  Changing the representative of
$Q$ by $A_R$ does not change the displayed coset, since
$R\subseteq T\setminus E(Q)$.

We verify order preservation with representatives included.  If
$Q_1=A_{R_1}z_1\subseteq Q_2=A_{R_2}z_2$, then all colors outside
$R_2$ have the same actual data at $z_1$ and $z_2$.  Hence
\[
 E(Q_2)=E(Q_1)\cap(T\setminus R_2).
\]
Putting $K_a=T\setminus E(Q_a)$ gives $K_1\subseteq K_2$ and
$R_2\subseteq K_2$.  Moreover
$z_1\in Q_2\subseteq A_{K_2}z_2$.  Thus
$A_{K_1}z_1\subseteq A_{K_2}z_2$, as required.  Finally, if every open
color of $Q$ is eligible, then $T\setminus E(Q)=R$ and $\rho(Q)=Q$.
So \eqref{eq:global-closure} is the stated order retraction.
The order homotopy observation in Section~\ref{sec:residues}
identifies the homotopy type of the earlier lower part with that of
its all-eligible image.
\end{proof}

\begin{lemma}\label{lem:global-residue-blocks}
The all-eligible lower poset at $P=A_Tx$ is a product of the residue
sublevel posets of Proposition~\ref{prop:residue-sublevels}, each with a
largest element adjoined, with the all-largest tuple deleted.  More
precisely, its factors are as follows.
\begin{enumerate}
 \item Each connected component $B$ of $T$ with negative sign at its
 highest level $q_B$ contributes the strict sublevel in the coset
 $A_Bx$ at level $q_B$.
 \item For a positive sign, group all components of $T$ belonging to the
 same positive-support component $I$ at the same level $q$.  Their union
 is $B=T\cap I$, and it contributes the special closed sublevel in
 $A_Bx$ with data $I,q,c_I=p_I^q(P)$.
\end{enumerate}
Every strict factor is empty or contractible, and every special closed
factor is nonempty and contractible.
\end{lemma}

\begin{proof}
Consider first a connected component $B$ with sign $-1$ and highest
level $q_B$.  Its boundary is outside $T$, because $B$ is an entire
component of $T$.  It has a cap at $q_B$: if some exterior neighbor had
a positive coordinate at a level $d>q_B$, the harmonic equation on $B$
and positivity of $C_B^{-1}$ would give a positive coordinate throughout
$B$ at that level, contrary to its definition.  On this boundary the
actual data $m^d(x)$ equal $p^d(P)$.  Their harmonic extension on $B$
therefore equals $p_B^{q_B}(P)$ at the test level.  Eligibility is
exactly a cap followed by the strict inequality against this harmonic
vector.  These are all the hypotheses of the strict residue sublevel.

Now take a positive-sign component, with support component $I$ at level
$q$.  Condition~\eqref{eq:global-null-sign} gives $q>0$, a positive
vector $c_I=p_I^q(P)$, and vanishing of all higher coordinates on
$I\cup\partial I$.  The coordinates on $\partial I$ at level $q$ are
zero by the definition of a support component.  Hence
$(Cp^q(P))_I=0$ also gives $C_Ic_I=0$.

Every component of $T$ meeting $I$ lies wholly in $I$, by harmonic
positivity on that component at level $q$.  Its highest level is $q$
because the higher coordinates vanish on $I$, and it has the same
positive sign.  We must group all these components, giving $B=T\cap I$.
If a vertex $j\in T\setminus I$ neighbored $I$, then $p_j^q(P)=0$,
since otherwise it would lie in the same support component.  An edge
to a positive coordinate in $I$ would give $(Cp^q(P))_j<0$, contradicting
harmonicity on $T$.  Therefore
\begin{equation}\label{eq:global-closed-exterior}
 \partial I\cap T=\varnothing,
 \qquad (I\setminus B)\cup\partial I
           \subseteq\widehat S\setminus T.
\end{equation}
All exterior values required by the special closed sublevel are thus
actual, fixed values on the entire coset $A_Tx$.  They are capped at
$q$ and equal the zero-extended vector $c$ at that level.  The block
$B$ is spherical and its harmonic prescription is $c_B$.
It cannot equal $I$, because a positive-definite principal matrix cannot
have the positive null vector $c_I$.  Eligibility on $B$ is exactly the
closed inequality in Proposition~\ref{prop:residue-sublevels}.  This
verifies every special closed hypothesis, including when several
disconnected components of $T$ meet the same $I$.

These blocks partition $T$ into unions of its connected components.
The presentation and the standard embeddings give
$A_T=\prod_B A_B$.  For $z\in A_Tx$, write uniquely
$zx^{-1}=\prod_B g_B$ with $g_B\in A_B$.  A subcoset $A_Rz$ corresponds
to the tuple
\[
 \bigl(A_{R\cap B}g_Bx\bigr)_B.
\]
Coset equality and inclusion are coordinatewise under this
correspondence.  Operations in other blocks leave every multiplicity in
$B$ unchanged, so the eligible-color conditions in $B$ are precisely
those for $g_Bx$ in the displayed block coset.  The boundary data just
verified lie outside $T$: for a strict block this follows because it is
a full component of $T$, and for a closed block it is
\eqref{eq:global-closed-exterior}.  They are therefore fixed under every
operation in $A_T$, including operations in the other blocks.  Opening nothing in a block gives its
single largest element $A_Bx$; opening somewhere gives the stated
proper residue sublevel.  The requirement $R\subsetneq T$ deletes just
the all-largest tuple.  This proves the product description.  Its
contractibility assertions follow from
Proposition~\ref{prop:residue-sublevels}.
\end{proof}

\begin{lemma}\label{lem:global-product-minus-top}
Let $\mathcal E_1,\ldots,\mathcal E_r$ be posets whose realizations are
each empty or contractible, and let $\mathcal E_a^+$ denote the result
of adjoining a largest element.  Then
\[
 \left(\prod_{a=1}^r\mathcal E_a^+\right)
          \setminus\{(\top,\ldots,\top)\}
\]
is empty if every $\mathcal E_a$ is empty, and otherwise has
contractible realization.  The empty product has the same empty
conclusion after its largest element is deleted.
\end{lemma}

\begin{proof}
We first record the product fact in the topology of order complexes.
If $\mathcal P$ and $\mathcal Q$ are nonempty posets with contractible
realizations, projection $\mathcal P\times\mathcal Q\to\mathcal P$
has lower fiber $\mathcal P_{\le p}\times\mathcal Q$ over $p$.
The order map $(p',q)\mapsto(p,q)$ is pointwise above the identity and
retracts that fiber onto $\{p\}\times\mathcal Q$.  The fiber is
contractible.  Quillen's Theorem~A, in the form recalled in
Section~\ref{sec:topology}, shows that the product poset has contractible
realization.  Iteration proves this for any finite number of nonempty
contractible factors.  This argument does not require identifying an
infinite order-complex realization with an ordinary topological product.

Discard indices for which $\mathcal E_a$ is empty, since their augmented
factor has only its largest element.  If no indices remain, the deleted
product is empty.  Otherwise cover its realization by the subcomplexes
$X_a$ consisting of chains whose $a$th coordinate lies in
$\mathcal E_a$.  They cover every chain: its largest vertex has a
nonlargest coordinate, and every earlier vertex has a nonlargest
coordinate at that same index.  For a nonempty set $J$ of indices,
the intersection $\bigcap_{a\in J}X_a$ is the realization of
\[
 \left(\prod_{a\in J}\mathcal E_a\right)
        \times
 \left(\prod_{a\notin J}\mathcal E_a^+\right).
\]
It is nonempty and contractible by the product fact: the required
factors are contractible, and each augmented factor has a largest
element.

There are finitely many subcomplexes in this cover.  Their union is
contractible by induction on their number.  Indeed the intersection of
the last subcomplex with the preceding union has a cover by the
corresponding intersections; all of its nonempty finite intersections
are again nonempty contractible.  Induction makes both the preceding
union and this intersection contractible.  Gluing the last
contractible subcomplex along the contractible intersection preserves
the homotopy type of a point, since subcomplex inclusions are
cofibrations.  This completes the proof.
\end{proof}

\begin{proposition}\label{prop:global-lower-link}
The earlier strict lower part of every $P\in\mathcal D$ is empty or
contractible.
\end{proposition}

\begin{proof}
For type $T=\varnothing$ the lower part is empty.  Otherwise
Lemma~\ref{lem:global-lower-comparison} retracts it onto the all-eligible
part.  Lemma~\ref{lem:global-residue-blocks} identifies that part with
the product in Lemma~\ref{lem:global-product-minus-top}, whose factors
satisfy that lemma's hypotheses.  The conclusion follows.
\end{proof}

\subsection{Contractibility and the Salvetti complex}

\begin{theorem}\label{thm:deligne-contractible}
For every Coxeter matrix on a finite set $S$, the order complex of the
spherical-coset poset $\mathcal D$ is contractible.
\end{theorem}

\begin{proof}
If $S=\varnothing$, the poset has a single vertex.  Otherwise use the
well-order of Lemma~\ref{lem:global-well-order}.
Propositions~\ref{prop:global-upper-link} and
\ref{prop:global-lower-link} show that the earlier strict upper and
strict lower parts at each vertex are empty or contractible.  Their
join is the earlier attaching link.  The empty-part convention and
Lemma~\ref{lem:filtration-principle} therefore show that each component
of the final order complex is contractible.

Finally $\mathcal D$ is connected.  The singleton cosets $x$ and
$\sigma_sx$ both lie in $A_{\{s\}}x$, since rank-one types are
spherical.  Left generator steps therefore connect all singleton
cosets.  Every other coset contains a singleton coset and hence joins
this same component.  The preceding componentwise conclusion proves
contractibility of $|\mathcal D|$.
\end{proof}

\begin{proof}[Proof of Theorem~\ref{thm:main}]
Lemma~\ref{lem:salvetti-comparison} gives a homotopy equivalence from
the universal cover of the full Artin Salvetti complex to
$|\mathcal D|$.  The latter is contractible by
Theorem~\ref{thm:deligne-contractible}, so the universal cover is
contractible.  Since covering maps induce isomorphisms on homotopy
groups in degrees at least two, the specified Salvetti complex is
aspherical.
\end{proof}

\subsection{Spherical, affine, and product cases}

If $S$ itself is spherical, the coset $A_S$ is a largest element of
$\mathcal D$.  Its contractibility is therefore immediate and agrees
with the spherical theorem used in the comparison.  The proof for an
arbitrary $S$ uses finite Coxeter length, finite Artin monoids, and
positive-definite inverses only on spherical subsets and their finite
unfoldings.

The positive-null-vector alternative, including the affine case, is
handled by the reversed preference in \eqref{eq:global-null-sign}.

If the original diagram splits as $S=S_1\sqcup\cdots\sqcup S_t$ with
all cross labels equal to $2$, its presentation gives
$A=\prod_a A_{S_a}$, and a type $T$ is spherical exactly when every
$T\cap S_a$ is spherical.  The finite Coxeter permutahedra split as
products on these components.  Shortest representatives and their
positive lifts split as well, because lengths add and letters in
different factors commute.  Consequently the canonical Salvetti cell
and face maps give the product of the factor Salvetti complexes, and
their universal covers give the corresponding product of covers.  The
argument above is compatible with this decomposition; in particular
the lower-block construction allows disconnected spherical types.

An infinite label contributes weight $2$ and an edge object consisting
of two copies of the tensor unit, and imposes no Artin braid relation.  It is
excluded from a spherical subset by positive definiteness itself: the
two-color principal matrix with that label is singular.  No restriction
on infinite labels in the original matrix, or on the connectedness of
its diagram, was used in the contraction.
